\chapter{Phase 4 --- Real-Time FPGA Prototype on the PYNQ-Z2}\label{ch:fpga}

\weeks{12--16} \quad Milestone M3

Simulation proves the design is correct; the FPGA proves it is \emph{real}. In
this phase your RTL processes live audio in real time, and typed complaints
change what you hear by writing hardware registers. That is your thesis
architecture on a single chip, with real numbers for latency, resources and
power.

\begin{whybox}
Phase~4 turns RQ2 and RQ3 from estimates into measurements: end-to-end latency
through real converters, resource use on a real device, and measured cycles per
frame (H4). It is also the demo: a 60-second video of the system working is the
single most persuasive artefact in your portfolio.
\end{whybox}

\begin{center}\small
\begin{tabularx}{\textwidth}{@{}LL@{}}
\toprule
\textbf{Inputs} & \textbf{Outputs} \\
\midrule
M2: bit-exact, verified RTL; register map (\cref{app:registers}); a Vivado
host (\cref{ch:setup}); the hardware listed below
& Scripted Vivado project in \file{fpga/pynq_z2/}; bitstream and overlay; a
Python driver; a hardware-in-the-loop bit-exact test; latency, resource and
power measurements; demo video \\
\bottomrule
\end{tabularx}
\end{center}

% =========================================================================
\section{Hardware and host}\label{sec:fpga:hw}

\begin{itemize}
  \item \textbf{TUL PYNQ-Z2 board.} A Zynq XC7Z020 SoC (a dual-core ARM
  Cortex-A9 processing system plus FPGA fabric), 512\,MB DDR3, and an Analog
  Devices ADAU1761 audio codec with headphone and line/microphone jacks. It is
  supported by the PYNQ framework~\cite{pynq,ch12-13:pynqdocs}, which runs
  Linux and Jupyter on the ARM and loads FPGA designs (\enquote{overlays}) from
  Python.
  \item A microSD card (16\,GB or more), a USB cable, an Ethernet cable,
        headphones, and a line-level audio source.
  \item Optional: your thesis's RME Fireface for loopback latency measurements,
        and a USB power meter.
  \item \textbf{Host:} Vivado on an x86-64 Linux or Windows machine
        (\cref{ch:setup}). Use the Vivado version that your PYNQ image release
        names in its documentation, so the board files and drivers match.
\end{itemize}

\begin{conceptbox}{Processing system and programmable logic}
A Zynq has two halves. The \emph{processing system} (PS) is a dual-core ARM
running Linux, with the DDR memory. The \emph{programmable logic} (PL) is FPGA
fabric where your RTL lives. They communicate over AXI buses:
\begin{itemize}
  \item AXI4-Lite for registers, so your \code{regs\_axil} appears as a block of
        memory addresses to Python;
  \item AXI4-Stream plus a DMA engine for bulk data.
\end{itemize}
This split matches your architecture: control plane on the PS, data plane in
the PL~\cite{ch12-13:crockett2014zynq}.
\end{conceptbox}

% =========================================================================
\section{Steps}\label{sec:fpga:steps}

\begin{stepbox}[label=st:fpga:boot]{Bring up the board and prove the audio path}
Flash the PYNQ image for the PYNQ-Z2 to the microSD card, boot, and open
Jupyter from your browser. Run the base overlay's audio example: record a few
seconds and play them back through the headphones.

This proves the codec, cables and jacks before any of your own hardware is
involved. Read how the base overlay's Python audio driver configures the
ADAU1761. You will reuse that configuration sequence.

\textbf{Check:} you hear your own recording. Note the PYNQ image version and
the Vivado version it requires in your notebook. \textbf{Effort:} half a day.
\end{stepbox}

\begin{stepbox}[label=st:fpga:dma]{Hardware-in-the-loop first: \texttt{ise\_top} behind an AXI DMA}
Before touching real-time audio, get your design on the chip in the simplest
possible system:
\begin{itemize}
  \item the Zynq PS;
  \item an AXI DMA engine that streams samples from DDR into \code{ise\_top} and
        streams results back;
  \item \code{regs\_axil} on AXI4-Lite.
\end{itemize}
Wrap \code{ise\_top} in a thin adapter with AXI4-Stream input and output, and
package it as IP.

Create the whole Vivado project from a Tcl script in
\file{fpga/pynq_z2/build.tcl}, never by clicking only, so it is version
controlled and reproducible. A good habit is to build the block design once in
the GUI, then export it with \code{write\_bd\_tcl} and commit the script.

\begin{lstlisting}[style=py,caption={Hardware-in-the-loop test from a PYNQ notebook (API names as in current PYNQ; check your version).}]
import numpy as np
from pynq import Overlay, allocate

ol = Overlay("ise_hil.bit")                     # loads bitstream + .hwh
dma = ol.axi_dma_0
x = np.load("golden/vectors/utt01_input.npy").astype(np.int16)
y_ref = np.load("golden/vectors/utt01_output.npy").astype(np.int16)

src = allocate(shape=x.shape, dtype=np.int16); src[:] = x
dst = allocate(shape=y_ref.shape, dtype=np.int16)
dma.recvchannel.transfer(dst); dma.sendchannel.transfer(src)
dma.sendchannel.wait(); dma.recvchannel.wait()
print("bit-exact:", np.array_equal(np.asarray(dst), y_ref))
\end{lstlisting}

\textbf{Check:} the FPGA output is bit-exact with the golden model, the same
result as in simulation, now on silicon. \textbf{Effort:} 3--5 days. The first
Vivado project always takes longer than expected.
\end{stepbox}

\begin{stepbox}[label=st:fpga:timing]{Timing closure and utilisation}
Run implementation at your chosen PL clock (\qty{100}{\MHz} from the PS is a
convenient start), then read two reports:
\begin{itemize}
  \item \code{report\_timing\_summary}: worst negative slack (WNS) must be
        $\ge 0$;
  \item \code{report\_utilization}: LUTs, flip-flops, DSP slices and block RAM
        tiles, both per block and in total.
\end{itemize}
Compare block RAM use with your state-memory calculation from \cref{ch:workload}.

\textbf{Check:} timing is met with margin. Utilisation is recorded in
\file{docs/resources.md}. If timing fails, add one pipeline register on the
failing path. Your design only needs a few MHz of real throughput, so it does
not need to run at the top clock speed. \textbf{Effort:} 1 day.
\end{stepbox}

\begin{stepbox}[label=st:fpga:audio]{Real-time audio: I2S to the codec}
Now replace the DMA path with the codec. You need three things:
\begin{enumerate}
  \item \textbf{Codec configuration.} Send the ADAU1761 register sequence over
  I2C from the PS. Reuse what the base overlay's driver does, together with the
  codec datasheet, to set the sample rate (\qty{48}{\kHz}, or the rate chosen in
  Phase~1), the I2S format and the signal routing.
  \item \textbf{I2S receiver and transmitter} in the PL (\file{rtl/i2s_rx.sv},
  \file{rtl/i2s_tx.sv}). They deserialise and serialise samples using the bit
  clock and word-select (LRCLK) signals.
  \item \textbf{Clock-domain crossing.} The I2S clocks are not your system clock.
  Bring samples across through a small asynchronous FIFO, or capture each
  sample in the I2S domain and pass a synchronised strobe (a two-flip-flop
  synchroniser) with the data held stable. Never sample a multi-bit bus
  directly across clock domains.
\end{enumerate}

\textbf{Check:} with \code{CTRL.BYPASS} set, audio passes through unaltered,
apart from delay. With \code{BYPASS} clear and default parameters, you hear
noise reduction. With $\beta$ raised through the registers, you hear stronger
suppression, as in the thesis. \textbf{Effort:} 1--1.5 weeks.
\end{stepbox}

\begin{stepbox}[label=st:fpga:control]{Close the loop: complaints control the hardware}
Write a small Python service on the PYNQ that receives parameter updates and
applies them with the driver in \cref{app:registers}. Then choose where the NLP
runs:
\begin{description}[style=nextline,leftmargin=1.2em]
  \item[Option A (recommended first): NLP on your laptop]
  Run the thesis pipeline unchanged: mpnet, SVM and Ridge, and the parameter
  mapping. Send the resulting parameters to the board over the network. This
  mirrors a real product architecture, where language processing runs on a
  paired phone and the hearing aid receives only parameters.
  \item[Option B (stretch): NLP on the board]
  Run a small encoder on the Cortex-A9. Be realistic: 512\,MB of RAM rules out
  fp32 mpnet, and many ML runtimes ship only 64-bit ARM builds while the
  Zynq-7000 is 32-bit. Check the availability of your runtime before you
  commit to this. If it works, measure the latency per sentence. This feeds
  directly into Track~A (\cref{ch:mlaccel}).
\end{description}

\textbf{Check:} typing \enquote{the background noise is unbearable} audibly
increases suppression within a fraction of a second, with no clicks at the
switch. \textbf{Effort:} 2--3 days.
\end{stepbox}

\begin{stepbox}[label=st:fpga:measure]{Measure latency, cycles and power}
\begin{description}[style=nextline,leftmargin=1.2em]
  \item[Latency]
  Use the thesis's RME Fireface. Play a chirp or click into the board's line
  input while recording both the source and the board's output on two Fireface
  inputs. The lag of the cross-correlation peak is the end-to-end latency. Then
  measure again in \code{BYPASS} mode. The difference is your algorithm and
  datapath; the bypass value is the converters and I2S buffering. Compare with
  the algorithmic latency predicted for your frame configuration.
  \item[Cycles]
  Read \reg{CYCLES\_LAST} over many frames. Its maximum, against the frame
  budget, is your measured answer to H4.
  \item[Power]
  Report Vivado's \code{report\_power} estimate for the PL (static and dynamic),
  ideally using switching activity from a simulation of real audio rather than
  the default activity rates. Then measure board power with a USB meter, with
  \code{CTRL.ENABLE} on and off. Expect the board total to be dominated by the
  ARM, the DDR memory and the regulators. That is exactly why the Vivado PL
  estimate and energy per frame, $E_\text{frame}=P_\text{dyn}\,T_\text{frame}$,
  are the meaningful numbers.
\end{description}

\begin{center}\footnotesize
\begin{tabularx}{\textwidth}{@{}LCCC@{}}
\toprule
Measurement & Value & Method & Notes \\
\midrule
End-to-end latency (processing) & & loopback cross-correlation & \\
End-to-end latency (bypass) & & loopback cross-correlation & \\
Max cycles per frame / budget & & \reg{CYCLES\_LAST} & H4 \\
LUT / FF / DSP / BRAM & & \code{report\_utilization} & \\
PL dynamic / static power & & \code{report\_power} & activity source? \\
Board power, enabled / disabled & & USB meter & \\
Energy per frame & & $P_\text{dyn}T_\text{frame}$ & \\
\bottomrule
\end{tabularx}
\end{center}

\textbf{Check:} the table is filled in, every value has its method and
conditions in the notebook, and the latency matches your prediction within a
few milliseconds of converter delay. \textbf{Effort:} 2--3 days.
\end{stepbox}

\begin{stepbox}[label=st:fpga:demo]{Record the demo}
Script a 60--90 second video:
\begin{enumerate}
  \item noisy speech playing;
  \item type a complaint; the parameters change;
  \item the audio audibly improves;
  \item type \enquote{that's perfect now} (\code{BALANCED}) and the parameters
        decay back towards their defaults;
  \item finish on the resource and latency numbers on screen.
\end{enumerate}
Record with good audio: use the Fireface, not a phone microphone next to
headphones.

\textbf{Check:} someone who has never seen the project understands it from the
video alone. \textbf{Effort:} 1 day.
\end{stepbox}

% =========================================================================
\section{Pitfalls}\label{sec:fpga:pitfalls}

\begin{pitfallbox}[Pitfalls in Phase 4]
\begin{itemize}
  \item \textbf{Vivado and PYNQ versions that do not match.} Board files and
        overlay metadata differ between releases. Use the documented pairing.
  \item \textbf{Starting with real-time audio.} Get the bit-exact DMA test
        working first. Then any problem in real time is in the I2S, clocking or
        codec setup, not in your datapath.
  \item \textbf{Clock-domain crossing by accident.} Every signal crossing
        between the I2S and system clocks needs a deliberate synchroniser or
        FIFO. Such bugs appear intermittently and are miserable to find.
  \item \textbf{Codec configuration.} A wrong I2S word length or a missed
        routing register gives silence or noise. Compare your register
        sequence with the working base overlay.
  \item \textbf{Address map surprises.} Check the address Vivado assigned to
        \code{regs\_axil} and read \reg{VERSION} first as a sanity check.
  \item \textbf{Comparing board power to a hearing aid.} The PYNQ board is not
        a low-power platform. Report PL estimates and energy per frame, and
        explain the difference (\cref{ch:evaluation}).
\end{itemize}
\end{pitfallbox}

\begin{readbox}
PYNQ documentation~\cite{ch12-13:pynqdocs} and its support
forum~\cite{ch12-13:pynqforum}. \emph{The Zynq Book}, for the PS/PL
architecture and AXI~\cite{ch12-13:crockett2014zynq}. FPGA DSP arithmetic and
resource mapping: Meyer-Baese~\cite{meyerbaese2014fpgadsp}.
\end{readbox}

\begin{donebox}[Milestone M3 is reached when]
\begin{checklist}
  \item The Vivado project is rebuilt from a committed Tcl script.
  \item The hardware-in-the-loop test is bit-exact on the FPGA.
  \item Timing is met; utilisation is recorded.
  \item Live audio is enhanced in real time; bypass works; parameter changes are click-free.
  \item Typed complaints control the hardware (option A at least).
  \item The latency, cycle, resource and power table is filled in with methods.
  \item The demo video is recorded.
\end{checklist}
\end{donebox}
